% Part: applied-modal-logics
% Chapter: epistemic-logic
% Section: introduction

\documentclass[../../../include/open-logic-section]{subfiles}

\begin{document}

\olfileid[id]{aml}{el}{int}

\olsection{Pengantar}

Sebagaimana logika modal membahas \emph{proposisi modal} dan relasi perikutan
di antaranya, logika epistemik membahas \emph{proposisi epistemik} dan relasi
perikutan di antaranya. Alih-alih menafsirkan operator modal sebagai
representasi kemungkinan dan keniscayaan, konektif uner ditafsirkan secara
epistemik atau doksastik untuk memodelkan pengetahuan dan keyakinan. Sebagai
contoh, kita mungkin ingin menyatakan hal-hal berikut:

\begin{enumerate}
\item Richard tahu bahwa Calgary berada di Alberta.
\item Audrey menganggap mungkin bahwa seekor anjing berada di sofa.
\item Richard tahu bahwa Audrey tahu kelasnya berlangsung setiap hari Selasa.
\item Semua orang tahu bahwa satu tahun terdiri atas 12 bulan.
\end{enumerate}

Logika epistemik kontemporer sering ditelusuri kembali ke \emph{Knowledge and
Belief} karya Jaako Hintikka dari tahun 1962, yang ditulis ketika semantik
dunia mungkin semakin banyak digunakan dalam logika. Logika epistemik pada
kenyataannya menggunakan sebagian besar perangkat semantik yang sama dengan
logika modal lain, tetapi menafsirkannya secara berbeda. Perubahan utamanya
terletak pada apa yang direpresentasikan oleh \emph{relasi aksesibilitas}.
Dalam logika epistemik, relasi-relasi tersebut merepresentasikan suatu bentuk
\emph{kemungkinan epistemik}. Kita akan melihat bahwa gagasan epistemik yang
kita modelkan memengaruhi syarat yang ingin kita kenakan pada relasi
aksesibilitas. Kita juga akan melihat apa yang terjadi pada teori
korespondensi ketika teori itu diberi penafsiran epistemik. Perhatikan bahwa
contoh di atas menyebut dua agen, Richard dan Audrey, beserta hubungan antara
hal-hal yang diketahui oleh masing-masing agen. Logika epistemik yang akan
kita tinjau merupakan logika multiagen, tempat hal-hal semacam itu dapat
dinyatakan. Sebaliknya, logika epistemik agen tunggal hanya membahas apa yang
diketahui atau diyakini oleh satu individu.

\end{document}
